% GENERATED FILE. Edit canonical lessons, then run npm run generate:books.
% canonical-source-hash: fnv1a64:5552bd2f0a69c5ab
% canonical-lessons: TE-C383-vadana, TE-C383-saksyam, TE-C383-nirupincu, TE-C383-tarkam, TE-C383-sahetukamaina

\chapter{Making a Case}
\label{ch:te-making-a-case}
\hlchaptermodality{383}

\begin{chapteropening}
\textbf{By the end of this chapter:} \emph{I can talk about arguments, evidence, proof and logic.}

Complete the last lesson of chapter 383: I can talk about arguments, evidence, proof and logic.
\end{chapteropening}

\section[\texorpdfstring{\te{వాదన} (vādana)}{వాదన (vādana)}]{\te{వాదన} (vādana) — a line of reasoning; an argument}

\label{lesson:TE-C383-vadana}



\begin{quote}
Before the new one: say the Telugu for to regard as, then the Telugu for an angle.
\end{quote}

\subsection*{You'll want to know: \te{వాదన}}
\textbf{\te{వాదన}} — \emph{vādana} — \textquotedblleft{}a line of reasoning; an argument\textquotedblright{}.

\begin{quote}
\te{అతని} \te{వాదనలో} \te{బలం} \te{లేదు}.
\end{quote}

\begin{quote}
\emph{atani vādanalō balaṁ lēdu.}
\end{quote}

\begin{quote}
There is no strength in his argument.
\end{quote}

Plural: \te{వాదనలు}, \emph{vādanalu}.

The case you make, not a quarrel. Its verb is \te{వాదించు}, to argue, which you know.

\begin{cousinweb}
\textbf{\te{వాదన}} (\emph{vādana}) is built on Sanskrit \emph{vāda}, \textquotedblleft{}speech, statement\textquotedblright{} (CDIAL 11507), from the root \emph{vad}, \textquotedblleft{}to speak\textquotedblright{}; the verb \te{వాదించు} comes from the same word.
\end{cousinweb}

\begin{grammarlens}[title={What you've built}]
One more word for giving reasons and opinions.
\end{grammarlens}

\subsection*{Your turn}
\begin{itemize}
\raggedright
  \item \emph{Say it:} \emph{vādana}
  \item \emph{Say it:} \emph{vādana}, once more
  \item \emph{Say it:} \emph{atani vādanalō balaṁ lēdu.} — There is no strength in his argument.
  \item \emph{Recall:} say the Telugu for to regard as, then the Telugu for an angle, then say \emph{vādana} again
\end{itemize}

\begin{tcolorbox}[breakable,colback=teal!4,colframe=teal!35!black,title={Before you move on}]
What does \te{వాదన} mean? (\textquotedblleft{}A line of reasoning; an argument\textquotedblright{}.) Say it once more.
\end{tcolorbox}
\section[\texorpdfstring{\te{సాక్ష్యం} (sākṣyaṁ)}{సాక్ష్యం (sākṣyaṁ)}]{\te{సాక్ష్యం} (sākṣyaṁ) — evidence; testimony}

\label{lesson:TE-C383-saksyam}



\begin{quote}
Before the new one: say the Telugu for an angle, then the Telugu for a line of reasoning.
\end{quote}

\subsection*{You'll want to know: \te{సాక్ష్యం}}
\textbf{\te{సాక్ష్యం}} — \emph{sākṣyaṁ} — \textquotedblleft{}evidence; testimony\textquotedblright{}.

\begin{quote}
\te{ఆ} \te{మాటకు} \te{సాక్ష్యం} \te{ఏమిటి}?
\end{quote}

\begin{quote}
\emph{ā māṭaku sākṣyaṁ ēmiṭi?}
\end{quote}

\begin{quote}
What is the evidence for that statement?
\end{quote}

\te{సాక్ష్యం} \te{చెప్పు}, to give evidence, to testify.

\begin{cousinweb}
\textbf{\te{సాక్ష్యం}} (\emph{sākṣyaṁ}) is borrowed from Sanskrit \emph{sākṣya}, \textquotedblleft{}evidence\textquotedblright{} (CDIAL 13322), and also \textquotedblleft{}testimony, attestation\textquotedblright{} (Monier-Williams).
\end{cousinweb}

\begin{grammarlens}[title={What you've built}]
One more word for giving reasons and opinions.
\end{grammarlens}

\subsection*{Your turn}
\begin{itemize}
\raggedright
  \item \emph{Say it:} \emph{sākṣyaṁ}
  \item \emph{Say it:} \emph{sākṣyaṁ}, once more
  \item \emph{Say it:} \emph{ā māṭaku sākṣyaṁ ēmiṭi?} — What is the evidence for that statement?
  \item \emph{Recall:} say the Telugu for an angle, then the Telugu for a line of reasoning, then say \emph{sākṣyaṁ} again
\end{itemize}

\begin{tcolorbox}[breakable,colback=teal!4,colframe=teal!35!black,title={Before you move on}]
What does \te{సాక్ష్యం} mean? (\textquotedblleft{}Evidence; testimony\textquotedblright{}.) Say it once more.
\end{tcolorbox}
\section[\texorpdfstring{\te{నిరూపించు} (nirūpiñcu)}{నిరూపించు (nirūpiñcu)}]{\te{నిరూపించు} (nirūpiñcu) — to prove}

\label{lesson:TE-C383-nirupincu}



\begin{quote}
Before the new one: say the Telugu for a line of reasoning, then the Telugu for evidence.
\end{quote}

\subsection*{You'll want to know: \te{నిరూపించు}}
\textbf{\te{నిరూపించు}} — \emph{nirūpiñcu} — \textquotedblleft{}to prove\textquotedblright{}.

\begin{quote}
\te{ఆమె} \te{తన} \te{మాట} \te{నిజమని} \te{నిరూపించింది}.
\end{quote}

\begin{quote}
\emph{āme tana māṭa nijamani nirūpiñcindi.}
\end{quote}

\begin{quote}
She proved that what she said was true.
\end{quote}

Past (I …): \te{నిరూపించాను}, \emph{nirūpiñcānu}.

What you prove is often a clause ending in \te{అని}: ... \te{నిజమని} \te{నిరూపించు}, to prove that ... is true.

\begin{cousinweb}
\textbf{\te{నిరూపించు}} (\emph{nirūpiñcu}) is built on Sanskrit \emph{nirūpaṇa}, \textquotedblleft{}stating, determining, defining\textquotedblright{}, and also \textquotedblleft{}examination\textquotedblright{} (Monier-Williams), with the Telugu verb ending -\te{ఇంచు}.
\end{cousinweb}

\begin{grammarlens}[title={What you've built}]
One more word for giving reasons and opinions.
\end{grammarlens}

\subsection*{Your turn}
\begin{itemize}
\raggedright
  \item \emph{Say it:} \emph{nirūpiñcu}
  \item \emph{Say it:} \emph{nirūpiñcu}, once more
  \item \emph{Say it:} \emph{āme tana māṭa nijamani nirūpiñcindi.} — She proved that what she said was true.
  \item \emph{Recall:} say the Telugu for a line of reasoning, then the Telugu for evidence, then say \emph{nirūpiñcu} again
\end{itemize}

\begin{tcolorbox}[breakable,colback=teal!4,colframe=teal!35!black,title={Before you move on}]
What does \te{నిరూపించు} mean? (\textquotedblleft{}To prove\textquotedblright{}.) Say it once more.
\end{tcolorbox}
\section[\texorpdfstring{\te{తర్కం} (tarkaṁ)}{తర్కం (tarkaṁ)}]{\te{తర్కం} (tarkaṁ) — logic}

\label{lesson:TE-C383-tarkam}



\begin{quote}
Before the new one: say the Telugu for evidence, then the Telugu for to prove.
\end{quote}

\subsection*{You'll want to know: \te{తర్కం}}
\textbf{\te{తర్కం}} — \emph{tarkaṁ} — \textquotedblleft{}logic\textquotedblright{}.

\begin{quote}
\te{అతని} \te{మాటల్లో} \te{తర్కం} \te{లేదు}.
\end{quote}

\begin{quote}
\emph{atani māṭallō tarkaṁ lēdu.}
\end{quote}

\begin{quote}
There is no logic in what he says.
\end{quote}

Careful, step-by-step reasoning. The adjective is \te{తార్కిక}, logical.

\begin{cousinweb}
\textbf{\te{తర్కం}} (\emph{tarkaṁ}) is borrowed from Sanskrit \emph{tarka}, \textquotedblleft{}inquiry\textquotedblright{} (CDIAL 5714), which Monier-Williams also glosses \textquotedblleft{}reasoning, speculation\textquotedblright{}.
\end{cousinweb}

\begin{grammarlens}[title={What you've built}]
One more word for giving reasons and opinions.
\end{grammarlens}

\subsection*{Your turn}
\begin{itemize}
\raggedright
  \item \emph{Say it:} \emph{tarkaṁ}
  \item \emph{Say it:} \emph{tarkaṁ}, once more
  \item \emph{Say it:} \emph{atani māṭallō tarkaṁ lēdu.} — There is no logic in what he says.
  \item \emph{Recall:} say the Telugu for evidence, then the Telugu for to prove, then say \emph{tarkaṁ} again
\end{itemize}

\begin{tcolorbox}[breakable,colback=teal!4,colframe=teal!35!black,title={Before you move on}]
What does \te{తర్కం} mean? (\textquotedblleft{}Logic\textquotedblright{}.) Say it once more.
\end{tcolorbox}
\section[\texorpdfstring{\te{సహేతుకమైన} (sahētukamaina)}{సహేతుకమైన (sahētukamaina)}]{\te{సహేతుకమైన} (sahētukamaina) — reasonable; well-grounded}

\label{lesson:TE-C383-sahetukamaina}



\begin{quote}
Before the new one: say the Telugu for to prove, then the Telugu for logic.
\end{quote}

\subsection*{You'll want to know: \te{సహేతుకమైన}}
\textbf{\te{సహేతుకమైన}} — \emph{sahētukamaina} — \textquotedblleft{}reasonable; well-grounded\textquotedblright{}.

\begin{quote}
\te{అది} \te{సహేతుకమైన} \te{ప్రశ్న}.
\end{quote}

\begin{quote}
\emph{adi sahētukamaina praśna.}
\end{quote}

\begin{quote}
That is a reasonable question.
\end{quote}

A reason stands behind it: \te{స}-, with, + \te{హేతు}, a reason. The ending -\te{మైన} makes an adjective, as in \te{కఠినమైన}.

\begin{cousinweb}
\textbf{\te{సహేతుకమైన}} (\emph{sahētukamaina}) is built on Sanskrit \emph{sahetuka}, \textquotedblleft{}together with a reason\textquotedblright{} (Monier-Williams), from \emph{sa-}, \textquotedblleft{}with\textquotedblright{}, and \emph{hetu}, \textquotedblleft{}a cause, a reason\textquotedblright{}; Telugu adds -\te{మైన} to make an adjective.
\end{cousinweb}

\begin{grammarlens}[title={What you've built}]
One more word for giving reasons and opinions.
\end{grammarlens}

\subsection*{Your turn}
\begin{itemize}
\raggedright
  \item \emph{Say it:} \emph{sahētukamaina}
  \item \emph{Say it:} \emph{sahētukamaina}, once more
  \item \emph{Say it:} \emph{adi sahētukamaina praśna.} — That is a reasonable question.
  \item \emph{Recall:} say the Telugu for to prove, then the Telugu for logic, then say \emph{sahētukamaina} again
\end{itemize}

\begin{tcolorbox}[breakable,colback=teal!4,colframe=teal!35!black,title={Before you move on}]
What does \te{సహేతుకమైన} mean? (\textquotedblleft{}Reasonable; well-grounded\textquotedblright{}.) Say it once more.
\end{tcolorbox}
